\section{Introduction}\label{sec:intro}

A \emph{convex disk} is a compact convex subset of the plane with nonempty interior. For a convex
disk $K$ let $\delta(K)$ be its greatest packing density: the supremum of the densities of the
packings of the plane by congruent copies of $K$ (\cref{def:packing-density}). The circular disk
has greatest packing density $\pi/\sqrt{12}=0.906899\ldots$. A disk is centrally symmetric if
it is invariant under reflection through a point. In 1934 Reinhardt~\cite{Reinhardt1934} asked
which centrally symmetric convex disk minimizes $\delta(K)$. He proposed the \emph{smoothed octagon} $\Koct$: the regular octagon with each corner replaced by an
arc of the hyperbola tangent to the two incident sides and asymptotic to the two
neighbouring side lines (\cref{def:smoothed-octagon}, \cref{fig:octagon}). Its greatest packing
density is
\[
  \doct=\frac{8-4\sqrt2-\log2}{2\sqrt2-1}=0.902414\ldots .
\]
We prove that the smoothed octagon minimizes the greatest packing density, uniquely up to
invertible affine maps. Such a map has the form $x\mapsto Lx+b$, where $L$ is an invertible
linear map and $b$ is a translation vector.

\begin{theorem}\label{thm:main}
\begin{enumerate}[label=\textup{(\alph*)}]
\item The smoothed octagon $\Koct$ is a centrally symmetric convex disk with $\delta(\Koct)=\doct$.
\item Every centrally symmetric convex disk $K$ satisfies $\delta(K)\ge\doct$.
\item A centrally symmetric convex disk $K$ satisfies $\delta(K)=\doct$ if and only if
  $K=A(\Koct)$ for an invertible affine map $A$ of the plane.
\end{enumerate}
\end{theorem}

Part~(b) is Reinhardt's conjecture. Part~(a), due to Reinhardt and
Mahler~\cite{Reinhardt1934,Mahler1947c}, is proved here in the form needed for the subsequent
argument. Part~(c) gives the uniqueness statement. Since $\delta$ is invariant under invertible
affine maps, its \emph{if} direction follows from~(a).

\subsection{History}\label{sec:intro-history}

Minkowski~\cite{Minkowski1904,Minkowski1907} studied the densest lattice packings of a convex
body and described the relevant lattices for a convex disk through inscribed hexagons. The
question of the least dense centrally symmetric disk appears in Blaschke's
book~\cite{Blaschke1923}. According to Hales and Vajjha~\cite[\S1.2]{HalesVajjha2024}, Blaschke
credits it to Courant, who expected the circle to be the answer. Reinhardt~\cite{Reinhardt1934}
proved that a least dense disk exists and has no corners. He also showed that every boundary
point is the midpoint of a side of a circumscribed hexagon of least area. His smoothed octagon
packs worse than the circle, but he left open whether it is the least dense disk.

Mahler, unaware of Reinhardt's work, reached the same conjecture in three papers of
1947~\cite{Mahler1947a,Mahler1947b,Mahler1947c}. He reduced the problem to one in the calculus of
variations and proved that a solution exists and is not an ellipse. He expected the boundary of an
extreme disk to consist of line segments and hyperbolic arcs~\cite[p.~117]{Mahler1947b}. Hales and
Vajjha call this Mahler's first conjecture. Mahler also showed that the regular octagon minimizes
density among centrally symmetric octagons. He constructed the smoothed octagon and computed its
density~\cite[\S12]{Mahler1947c}. Ledermann and Mahler~\cite{LedermannMahler1949} named the
smoothed octagon, and found that the smoothed decagon packs better, ``a result which seems to
support Reinhardt's conjecture''.

These authors studied lattice packings. Fejes T\'oth~\cite{FejesToth1950} and
Rogers~\cite{Rogers1951} proved that, for a centrally symmetric convex disk, packings by congruent
copies and packings by translates have the same greatest density as lattice packings
(\cref{thm:fejes-toth}). Reinhardt's lattice problem therefore extends to all packings; see the
surveys of G.~Fejes T\'oth~\cite{GFejesToth2017} and of Bezdek and
Kuperberg~\cite{BezdekKuperberg2013}. Lower bounds for $\delta(K)$ followed:
$\sqrt3/2=0.866\ldots$ by Mahler~\cite{Mahler1946}, $0.88$ by Ennola~\cite{Ennola1961}, and
$0.892656$ by Tammela~\cite{Tammela1970}. Nazarov~\cite{Nazarov1986} proved that the smoothed
octagon is a local minimum of $\delta$ in the Hausdorff metric.

Hales recast the problem twice. In~\cite{Hales2011} he formulates it as a variational problem
with an integral cost and prescribed boundary conditions. The corresponding Euler--Lagrange
equations have only the circle as a solution, up to affine maps, but the circle is not a
minimizer. He also proves that a piecewise analytic minimizer is a smoothed polygon.
In~\cite{Hales2017} it is an optimal control problem: the motion of six boundary contact points
is governed by differential equations, and the area of the disk is the quantity to be minimized.
Hales identifies the motions associated with smoothed regular $(6k+2)$-gons, proves that the
octagon is a strict local minimum, and gives conditions under which an optimal motion has only
finitely many changes of control. Vajjha's thesis~\cite{Vajjha2022} and the book of Hales and
Vajjha~\cite{HalesVajjha2024} continue this work. The book confines the state of a global
minimizer to a compact domain (\cref{thm:compact-domain}) and proves Mahler's first conjecture:
a worst disk is a smoothed polygon, with finitely many straight segments and hyperbolic arcs.
It leaves the Reinhardt conjecture open, and notes that a classification of the possible optimal
motions by the smoothed $(6k+2)$-gons would complete it~\cite[\S1.3]{HalesVajjha2024}; see
also~\cite{Barber2024}.

Related extremal packing problems have also been studied. In three dimensions,
Kallus~\cite{Kallus2014} proved that the ball is a local pessimum for packing. For convex disks
that need not be centrally symmetric, F\'ary~\cite{Fary1950} proved that triangles have the least
lattice packing density, and Kallus~\cite{Kallus2015} proved that the regular heptagon is a local
minimum of the double-lattice packing density.

\subsection{The strategy of the lower bound}\label{sec:strategy}

The first reduction concerns \emph{lattice packings}, in which the copies are translates of the
disk and their centres are the integer linear combinations of two fixed independent vectors. The
parallelogram spanned by those vectors is a \emph{lattice cell}. For a centrally symmetric convex
disk, the theorem of Fejes T\'oth and Rogers says that a densest packing can be chosen in this
form. Its density is the area of the disk divided by the cell area. An invertible linear map
changes both areas by the same factor. We may therefore normalize the cell area to $\sqrt{12}$ and
seek a lower bound for the area of the disk.

A \emph{worst disk} is a centrally symmetric convex disk that minimizes the greatest packing
density. Such a disk exists. Its densest lattice packings admit six boundary contact points that
move continuously around the disk while preserving the cell area. These points form an inscribed
hexagon and together trace the entire boundary. Hales's control problem describes this motion by
differential equations. The \emph{control} specifies the allowed curvatures of the contact curves,
and the \emph{cost} of the motion is the area of the traced disk. An \emph{admissible trajectory}
is a motion satisfying these equations and the closing conditions: after one period, each contact
point has reached the starting position of the next one. The precise formulation is in
\cref{def:control-problem}.

The smoothed octagon supplies an admissible trajectory of cost
\[
  \Aoct=\sqrt{12}\,\doct.
\]
The geometric problem is equivalent to the following bound.

\begin{theorem}\label{thm:control}
Every admissible trajectory of Hales's control problem has cost at least $\Aoct$.
\end{theorem}

The proof reduces the possible motions in three stages.
\begin{enumerate}[label=(\arabic*)]
\item \emph{A minimizer has finitely many pieces.}
A trajectory of least cost is a \emph{global minimizer}. Its control switches finitely often
among three distinguished choices, labelled $0$, $1$ and $2$. Each interval of constant control
is a \emph{piece}, and a change between consecutive pieces is a \emph{switch}. The equations
have explicit solutions on each piece, so the motion is determined by a finite sequence of
labels and durations.
\item \emph{Six or more pieces are impossible.}
At each switch the label advances or retreats in the cyclic order $0,1,2$. We call these switches
\emph{forward} and \emph{reverse}, respectively. A pattern with both directions is called
\emph{mixed}. The sequences fall into five classes: two with only one direction of switching
and three mixed classes. We exclude each class for a minimizer with six or more pieces. Some
patterns contradict local necessary conditions for optimality; the others force the cost of a
complete period above the octagon's. These exclusions form the main part of the proof.
\item \emph{The remaining motions can be compared exactly.}
With at most five pieces, the closing equations leave a four-piece family and a five-piece
family. The four-piece family has least cost $\Aoct$, attained only by the octagon's motion.
Every member of the five-piece family costs strictly more.
\end{enumerate}

For a worst disk $K_0$, the cost of its normalized motion is $\sqrt{12}\,\delta(K_0)$.
The comparison therefore gives
\begin{equation}\label{eq:proof-chain}
  \sqrt{12}\,\delta(K_0)\ge\Aoct=\sqrt{12}\,\doct.
\end{equation}
Every other centrally symmetric convex disk has density at least $\delta(K_0)$, which proves the
lower bound in \cref{thm:main}(b).

The exclusions relate local switching conditions to the cost of a complete period. The estimates
use monotonicity, convexity and integral bounds for the logarithm. Some mixed cases also require
polynomial sign checks, detailed in the corresponding appendix.

\subsection{The strategy of the equality case}\label{sec:strategy-equality}

For uniqueness, let $K$ be a centrally symmetric convex disk with $\delta(K)=\doct$. The lower
bound shows that, after translation of its centre to the origin, this disk is a worst disk. Its
trajectory is therefore a global minimizer of cost $\Aoct$.

The exclusions leave at most five pieces. The equality statement in the final comparison then
identifies this trajectory with the octagon's motion, up to a choice of starting point and a
cyclic permutation of the three control labels. Reversing these changes and the initial linear
normalization shows that the given disk $K$ is an invertible affine image of the smoothed octagon
(\cref{thm:rigidity}). The argument applies to the trajectory of $K$ itself and hence to every
disk attaining equality.

Conversely, invertible affine maps preserve the greatest packing density of centrally symmetric
convex disks (\cref{prop:converse}). Every affine image of the smoothed octagon consequently has
density $\doct$. This proves the equality characterization in \cref{thm:main}(c).

\subsection{Formal verification}\label{sec:intro-formal}

\Cref{thm:main,thm:control} are formalized in the proof assistant Lean~4~\cite{Lean4}, using its
mathematical library Mathlib~\cite{Mathlib20}. The proofs are available in the
repository~\cite{Companion}, and Lean's kernel checks every step. The three parts of
\cref{thm:main} are stated in Mathlib's vocabulary in a file that imports only Mathlib, and the
proof uses only Lean's standard axioms. It builds on formalizations of the book of Hales and
Vajjha~\cite{SmoothedPolygonLean} and of the Pontryagin maximum
principle~\cite{OptimalControlLean}, and it is registered in the Palomar
registry~\cite{Palomar}. \Cref{app:formal} describes the formal statements, their verification
and the corresponding Lean declarations.

\subsection{Organization}\label{sec:intro-organization}

\Cref{sec:setting} explains the passage from packing density to the control problem.
\Cref{sec:outline-periods} reduces a minimizer to finitely many pieces and classifies its
switches. \Cref{sec:outline-classes} gives the proof idea for each of the five exclusions.
\Cref{sec:outline-terminal} compares the remaining motions and completes the lower bound and
the equality case.

The appendices supply the details in the same order. \Cref{sec:reduction} gives the geometric
formulation and the reduction to finite normal periods. \Cref{sec:periods} develops the formulas
shared by the case arguments. Each of the five cases then has its own appendix:
\cref{sec:reverse,sec:mixed-runs,sec:short-runs,sec:forward,sec:mixed}. The two continuation
arguments within the last mixed case are subsections of \cref{sec:mixed}.
\Cref{sec:terminal,sec:equality} prove the final comparisons and recover the equality disks.
\Cref{app:formal} describes the formal verification. A notation table is provided in
\cref{sec:set-notation}.
